% Simple Beamer Template
\documentclass{beamer}
\beamertemplatenavigationsymbolsempty
\setbeamertemplate{footline}[page number]
\usepackage{amsmath}
\newcommand{\T}{^{\mathsf{T}}}
%\usepackage{array,tabularx,tabulary,booktabs}
%\usepackage{longtable}
%\usepackage{multirow}

\title{Causal Graphs and Optimization}
\author{IA pour la Science}
\date{February $16^\text{th}$, 2026}

\begin{document}
\begin{frame}
  \titlepage
\end{frame}
%=============================================
\begin{frame}{Questions to discuss}
\begin{enumerate}
	\item There is a difference between the causal graphs $G_t$ and $G_{t-1}$ for insufficient $T$.
    \item (Th.) There is no difference between the causal graphs $G_t$ and $G_{t-1}$ for sufficient $T$.
	\item The exhaustive and combinatorial algorithms are special cases of optimization.
    \item The latent space is a spectrum of the observable domain $t\to s$ as in Laplace transform.
    \item The time is continuous as in the Laplace transform with adaptation to the neural ODE\@.
\end{enumerate}
\end{frame}
%=============================================
\begin{frame}{MVP}
\begin{enumerate}
	\item Discrete time.
    \item Latent space.
	\item Direct and inverse causal graph.
    \item The target to optimize the causal graph.
    \item The graph structure.
    \item The natural gradient descent to optimize the model.
    \item The NGD estimate the distribution of the model parameters during Python optimization.
    \item Generative model to generate the inverse causal graph.
    \item Double pendulum example.
\end{enumerate}
\end{frame}
%=============================================
\begin{frame}{Covariance and Causality}
	\includegraphics[width=1.09\textwidth]{tmp/fig_discuss_dynamic_scm}
\end{frame}
%=============================================
\begin{frame}{Summary Causal Graph}
	\includegraphics[width=\textwidth]{tmp/fig_summary_causal_graph}
\end{frame}
%=============================================
\begin{frame}{The dimensionality of the model}
\[
    \begin{pmatrix}
        w_{1,1} & w_{1,2} & \cdots & w_{1,N} \\
        w_{2,1} & w_{2,2} & \cdots & w_{2,N} \\
        \vdots  & \vdots  & w_{\kappa,\tau} & \vdots  \\
        w_{K,1} & w_{K,2} & \cdots & w_{K,N}
    \end{pmatrix}
    \underbrace{
    \begin{pmatrix}
        x_{t-1} & x_{t-2} & \cdots & x_{t-(N+T)} \\
        x_{t-2} & x_{t-3} & \cdots & x_{t-(N+T-1)} \\
        \vdots & \vdots & x_{t-\tau} & \vdots \\
        x_{t-N} & x_{t-(N-1)} & \cdots & x_{T}
    \end{pmatrix}}_{\text{Hankel matrix }\mathbf{X}_{t-1}}.
\]
Convolution of $f$ and $g$
\[
	(f \ast g) (t) = \int_{\tau=0}^T f(\tau)g(t-\tau)d\tau
\]
The bilateral Laplace transform
\[
F(s) \cdot G(s) = \int_{-\infty}^\infty \int_{-\infty}^\infty e^{-st} \ f(u) \ g(t - u) \ \text{d}u \ \text{d}t
= \int_{-\infty}^\infty e^{-st} (f * g)(t) \ \text{d}t.
\]
\end{frame}
%=============================================
\begin{frame}{Phase Space of a pendulum: x,y-axis are $(\theta, \dot{\theta})$}

\centering{\includegraphics[width=0.65\textwidth]{tmp/fig_pendulum2_space}}

\[
    {\theta}(t) =-\mu\theta(t)  - \frac{g}{L} \sin\big(\theta(t)\big)
\]

\[
\frac{d}{dt}
\begin{bmatrix}
\theta(t) \\
\dot{\theta}(t)
\end{bmatrix}
=
\begin{bmatrix}
\dot{\theta}(t) \\
\ddot{\theta}(t)
\end{bmatrix}
=
\begin{bmatrix}
\dot{\theta}(t) \\
- \mu \dot{\theta}(t) - \frac{g}{L} \sin\big(\theta(t)\big)
\end{bmatrix}
\]

\end{frame}
%=============================================
\begin{frame}{Phase trajectory}
	\includegraphics[width=\textwidth]{tmp/fig_phase_trajectory}
\end{frame}
%=============================================
\begin{frame}{$\text{KL}\bigl(p(\mathbf{w}|X)\,||\,p(\mathbf{w}|do(X=\mathbf{x})\bigr)$}
	\includegraphics[width=\textwidth]{tmp/fig_discuss_wdox}
\end{frame}
%=============================================
    \begin{frame}{Graphical model for the structure parameters $\Gamma $}
	\includegraphics[width=0.5\textwidth]{tmp/fig_gamma}
\end{frame}

%=============================================
\begin{frame}{Causal graph optimization}

    \begin{columns}[T]
                \begin{column}{.5\textwidth}

    \begin{block}{There given}
    \begin{enumerate}
        \item Model structure $\Gamma$
        \item Causal graph $G$
        \item Causal hypergraph $\mathcal{G}$
        \item Observable $\mathbf{x}$ and latent $\mathbf{z}$ spaces and data with reconstructions
        \item Set of do-interventions in $\mathbf{x}$ and $\mathbf{z}$
    \end{enumerate}
    \end{block}

    \begin{block}{Set of criterions}
     ${X} \text{prind} {Y}$, $\text{KL}(X||Y)$, $\text{MI}(X,Y)$ and 3 rules of do-calculus
    \end{block}

            \end{column}
                \begin{column}{.5\textwidth}

    \begin{block}{The algorithm}
        Forward $(\mathbf{x} \to \mathbf{z})$ and backward $(\mathbf{z} \to \mathbf{x})$
    \begin{enumerate}
        \item Optimization step: get distributions (without do)
        \item Select a (hyper)node
        \item Run do-intervention
        \item Optimization step: get distributions (with do)
        \item Compute the criterion
        \item Eliminate the $v\in\mathcal{G}$ and $\gamma\in\Gamma$
        \item Continue until convergence
    \end{enumerate}
    \end{block}

                \end{column}
            \end{columns}

\end{frame}

%=============================================
\begin{frame}{Direct and indirect causal effects}
	\includegraphics[width=0.5\textwidth]{tmp/fig_causal_effect}


    Observe two-pendulum system with mass, or forced pendulum.
    The indirect causes are (are not) d-separated for
    $y$ is controlled by $z$, which has
    \begin{enumerate}
        \item no mass, no energy $\to$ are
        \item no mass, energy $\to$ are
        \item mass, no energy $\to$ are
        \item mass, energy $\to$ not
    \end{enumerate}
\end{frame}

\begin{frame}{Table left 1}
	\includegraphics[width=\textwidth]{../temporary/tab_IMG_4126_1.jpeg}
\end{frame}
%=============================================
\begin{frame}{Table left 2}
	\includegraphics[width=\textwidth]{../temporary/tab_IMG_4129_2.jpeg}
\end{frame}
%=============================================
\begin{frame}{Table left 3}
	\includegraphics[width=\textwidth]{../temporary/tab_IMG_4134_3.jpeg}
\end{frame}
    %=============================================
\begin{frame}{Table right 1}
	\includegraphics[width=\textwidth]{../temporary/tab_IMG_4131_4.jpeg}
\end{frame}
    %=============================================
\begin{frame}{Table right 12}
	\includegraphics[width=\textwidth]{../temporary/tab_IMG_4135_5.jpeg}
\end{frame}
\end{document}